% ============================================================================
% WORKBOOK — Section 1.5: Graphing Proportional Relationships
% CCSS 7.RP.A.2.D
% Practice-focused reinforcement for the study guide topic
% ============================================================================
\section{Graphing Proportional Relationships}
\topicTitle{Graphing Proportional Relationships}

\begin{quickReview}{Graphing Proportional Relationships}
A proportional relationship graphs as a \textbf{straight line through the origin} $(0,0)$.

\begin{itemize}[leftmargin=*]
    \item $(0, 0)$: zero of one quantity gives zero of the other.
    \item $(1, r)$: the $y$-value at $x = 1$ IS the \textbf{unit rate}.
    \item \textbf{Steeper line = greater unit rate.} This steepness is called the \textbf{slope}, and slope $= k$.
\end{itemize}

\medskip
\textbf{Example:} If a graph passes through $(0,0)$ and $(3, 12)$, then $k = 12 \div 3 = 4$.

At $x = 1$, the graph hits $(1, 4)$, so the unit rate is $4$.
\end{quickReview}

% ── Warm-Up ──────────────────────────────────────────────────────────────────
\begin{practiceBox}[funGreen]{\faSun~Warm-Up}

\practiceHeader[funGreen]{Find the unit rate from the given point on a proportional graph.}

\begin{multicols}{2}
\prob $(2, 10)$ \quad $k = $ \answerBlank[2cm]
\answerExplain{$k = 5$}{Divide $y$ by $x$: $10 \div 2 = 5$.}

\prob $(4, 18)$ \quad $k = $ \answerBlank[2cm]
\answerExplain{$k = 4.5$}{Divide $y$ by $x$: $18 \div 4 = 4.5$.}

\prob $(5, 35)$ \quad $k = $ \answerBlank[2cm]
\answerExplain{$k = 7$}{Divide $y$ by $x$: $35 \div 5 = 7$.}

\prob $(3, 7.5)$ \quad $k = $ \answerBlank[2cm]
\answerExplain{$k = 2.5$}{Divide $y$ by $x$: $7.5 \div 3 = 2.5$.}

\prob $(1, 6)$ \quad $k = $ \answerBlank[2cm]
\answerExplain{$k = 6$}{When $x = 1$, the $y$-value is the unit rate: $k = 6$.}

\prob $(8, 20)$ \quad $k = $ \answerBlank[2cm]
\answerExplain{$k = 2.5$}{Divide $y$ by $x$: $20 \div 8 = 2.5$.}
\end{multicols}
\end{practiceBox}

% ── Practice A: Interpreting Points ──────────────────────────────────────────
\begin{practiceBox}{What Do the Points Mean?}

\practiceHeader{Answer based on the proportional graph described.}

\prob The equation $y = 8x$ gives total cost of movie tickets. What does the point $(1, 8)$ mean in context?
\answerBlank[5cm]
\answerExplain{One movie ticket costs $\$8$}{The point $(1, r)$ gives the unit rate. Here $x = 1$ ticket and $y = \$8$, so one ticket costs $\$8$.}

\prob A graph of distance versus time passes through $(0, 0)$ and $(2, 90)$. What does $(0, 0)$ mean? What is the speed?
\answerBlank[5cm]
\answerExplain{$(0,0)$: no time means no distance; speed $= 45$ mi/hr}{The origin means $0$ hours traveled gives $0$ miles. The speed (unit rate) is $90 \div 2 = 45$ miles per hour.}

\prob A proportional graph passes through $(4, 6)$. What are the coordinates of the point at $x = 1$?
\answerBlank[3cm]
\answerExplain{$(1, 1.5)$}{Find $k$: $6 \div 4 = 1.5$. At $x = 1$, $y = 1.5$, so the point is $(1, 1.5)$.}

\prob Two proportional graphs share the point $(0, 0)$. Graph A passes through $(3, 15)$; Graph B passes through $(3, 9)$. Which graph is steeper?
\answerBlank[3cm]
\answerExplain{Graph A}{Graph A: $k = 15 \div 3 = 5$. Graph B: $k = 9 \div 3 = 3$. A higher $k$ means a steeper line, so Graph A is steeper.}

\end{practiceBox}

% ── Practice B: Comparing Lines ──────────────────────────────────────────────
\begin{practiceBox}[funTeal]{Comparing Proportional Graphs}

\practiceHeader[funTeal]{Compare the two relationships and answer.}

\prob Store A: $c = 3.50n$ \quad Store B: $c = 4.25n$ \quad Which store charges MORE per item? By how much?
\answerBlank[5cm]
\answerExplain{Store B charges $\$0.75$ more per item}{Store A: $k = 3.50$. Store B: $k = 4.25$. Store B charges more per item. Difference: $4.25 - 3.50 = \$0.75$.}

\prob Runner A passes through $(2, 12)$. Runner B passes through $(3, 15)$. Both lines start at $(0, 0)$. Who is faster?
\answerBlank[3cm]
\answerExplain{Runner A}{Runner A: $k = 12 \div 2 = 6$ units/time. Runner B: $k = 15 \div 3 = 5$. Since $6 > 5$, Runner A is faster.}

\prob Plant A grows $1.5$ cm per day. Plant B grows $2$ cm per day. After how many days will Plant B be $5$ cm taller than Plant A?
\answerBlank[3cm]
\answerExplain{$10$ days}{Difference per day: $2 - 1.5 = 0.5$ cm. To reach $5$ cm difference: $5 \div 0.5 = 10$ days.}

\prob Dani bikes at $12$ mph. Kim bikes at $8$ mph. They start at the same time. How far apart are they after $3$ hours?
\answerBlank[3cm]
\answerExplain{$12$ miles}{Dani: $12 \times 3 = 36$ miles. Kim: $8 \times 3 = 24$ miles. Difference: $36 - 24 = 12$ miles apart.}

\end{practiceBox}

% ── Practice C: True or False ────────────────────────────────────────────────
\begin{practiceBox}[funPurple]{True or False?}

\trueOrFalse{On a proportional graph, the point $(1, r)$ always shows the unit rate.}
\answerExplain{True}{By definition, $y = kx$ evaluated at $x = 1$ gives $y = k$. So the $y$-coordinate at $x = 1$ is the unit rate $k$.}

\trueOrFalse{A steeper proportional line always means a smaller unit rate.}
\answerExplain{False}{Steeper means a larger unit rate. A line with $k = 5$ is steeper than a line with $k = 2$ because $y$ grows faster per unit of $x$.}

\trueOrFalse{If two proportional graphs both pass through $(4, 20)$, they must be the same line.}
\answerExplain{True}{Both lines pass through $(0, 0)$ (proportional) and $(4, 20)$. Two points determine a unique line, so they are the same line with $k = 5$.}

\trueOrFalse{A proportional graph can cross the $y$-axis at a point other than $(0, 0)$.}
\answerExplain{False}{A proportional relationship always passes through the origin. If the $y$-intercept is anything other than $0$, the relationship is not proportional.}

\end{practiceBox}

% ── Practice D: Working from Equations ───────────────────────────────────────
\begin{practiceBox}[funBlue]{From Equations to Graphs}

\prob For $y = 6x$, find $y$ when $x = 0, 1, 2, 3$ and write the four points.
\answerBlank[5cm]
\answerExplain{$(0,0),\; (1,6),\; (2,12),\; (3,18)$}{Substitute: $6(0)=0$, $6(1)=6$, $6(2)=12$, $6(3)=18$. The four points are $(0,0)$, $(1,6)$, $(2,12)$, $(3,18)$.}

\prob A graph passes through $(5, 12.5)$ and $(0, 0)$. Write the equation and find $y$ when $x = 8$.
\answerBlank[4cm]
\answerExplain{$y = 2.5x$; $y = 20$}{Find $k$: $12.5 \div 5 = 2.5$. Equation: $y = 2.5x$. When $x = 8$: $y = 2.5 \times 8 = 20$.}

\prob A graph passes through $(0, 0)$ and $(4, 14)$. Find $x$ when $y = 49$.
\answerBlank[3cm]
\answerExplain{$x = 14$}{Find $k$: $14 \div 4 = 3.5$. Set $49 = 3.5x$: $x = 49 \div 3.5 = 14$.}

\end{practiceBox}

% ── Challenge ────────────────────────────────────────────────────────────────
\begin{challengeBox}
\prob Two proportional lines are graphed. Line A passes through $(2, 5)$. Line B passes through $(3, 5)$. Which line is steeper? Find the value of $x$ where Line B gives a $y$-value of $15$. 
\answerBlank[4cm]
\answerExplain{Line A is steeper; $x = 9$}{Line A: $k = 5 \div 2 = 2.5$. Line B: $k = 5 \div 3 \approx 1.67$. Since $2.5 > 1.67$, Line A is steeper. For Line B when $y = 15$: $15 = \frac{5}{3}x$, so $x = 15 \times \frac{3}{5} = 9$.}

\prob A proportional graph passes through $\left(\frac{1}{2}, 3\right)$. What is the slope? What is $y$ when $x = 4\frac{1}{2}$?
\answerBlank[4cm]
\answerExplain{Slope $= 6$; $y = 27$}{Find $k$: $3 \div \frac{1}{2} = 3 \times 2 = 6$. The slope is $6$. When $x = 4.5$: $y = 6 \times 4.5 = 27$.}
\end{challengeBox}

\encouragement{You can read proportional graphs like a map --- every point tells a story!}
